% arara: pdflatex
\documentclass{beamer}
\usepackage{my-proposal-talk}

\usepackage{lmodern}
\usepackage{fancyvrb}

\NewDocumentCommand{\fuse}{}{\bullet}
\NewDocumentCommand{\lrimp}{}{\twoheadrightarrow}
\NewDocumentCommand{\monad}{m}{\{#1\}}

\renewcommand*{\inc}{\color{structure}\mathtt{s}\color{black}}
\renewcommand*{\eps}{\color{structure}\mathtt{e}\color{black}}
\renewcommand*{\bit}[1]{\color{structure}\mathtt{#1}\color{black}}

\NewDocumentCommand \val { m }{ \llcorner #1 \lrcorner }
% \renewcommand*{\cntr}[1]{\ulcorner\color{example text.fg}#1\color{black}\urcorner}

\NewDocumentCommand \circtask { O{} m }{%
  \tikz { \fill [#1] (#2/2,#2/2) circle [radius={#2/2}]; }%
}
\NewDocumentCommand \sqtask { O{} m }{%
  \tikz { \fill [#1] (0,0) rectangle ++(#2,#2); }%
}
\NewDocumentCommand \tritask { O{} m }{%
  \tikz { \fill [#1]
            let \n1 = { 2*#2/sqrt(3) }
            in (0:0) -- ++(60:\n1) -- ++(-60:\n1) -- cycle; }%
}

\NewDocumentCommand \showcolumnwidth { }{%
  \tikz { \draw [|-|] (0,0) -- (\linewidth,0); }%
}

\usetikzlibrary{patterns}

\NewDocumentCommand \exec { m }{%
  \mathop{\mathsf{exec}} #1
}

\NewDocumentCommand \msg { m }{%
  \mathop{\mathsf{msg}} #1
}

\NewDocumentCommand \nats { m }{%
  \mathord{\!}\mathop{\mathsf{nats}}#1
}



% Presentation metadata
\title{Session-typed concurrent logic programming}
\author[DeYoung]{Henry DeYoung}
\institute[Carnegie Mellon]{%
  Computer Science Department\\Carnegie Mellon University}
\date{\today}
\titlegraphic{\textcolor{gray}{\rule{12em}{9em}}}



\begin{document}

\colorlet{exampletext}{example text.fg}

\begin{frame}[plain]
  \titlepage
\end{frame}


\begin{frame}[fragile]
  \frametitle{Computation and deduction}

  \vspace*{-8ex}

  \begingroup\Huge
    \begin{equation*}
      \Gamma \vdash A
    \end{equation*}%
  \endgroup

  \begin{columns}[T]
  \column{0.5\textwidth}
    \begin{tikzcd}[arrow style=tikz, /tikz/>={To[scale=0.8]}, arrows={very thick, start anchor=real east, end anchor=real west}, /tikz/nodes={inner xsep=0}]
      \begin{tikzpicture}[scale=0.75]
        \fill [gray] (0:0) -- ++(60:2) -- ++(0:-2) -- cycle;
        \clip (-60:-2.55) -- ++(-60:0.5) -- (60:2.05) -- ++(60:0.5) -- cycle;
        \fill [structure] (-60:-1.35) -- ++(60:1.2) -- ++(0:-1.2) -- cycle;
        \fill [example text.fg] (60:1.35) -- ++(60:1.2) -- ++(0:-1.2) -- cycle;
      \end{tikzpicture}
      \rar[structure] &
      \begin{tikzpicture}[scale=0.75]
        \fill [gray] (0:0) -- ++(60:2) -- ++(0:-2) -- cycle;
        \clip (-60:-2.55) -- ++(-60:0.5) -- (60:2.05) -- ++(60:0.5) -- cycle;
        \fill [example text.fg] (60:1.35) -- ++(60:1.2) -- ++(0:-1.2) -- cycle;
      \end{tikzpicture}
    \end{tikzcd}
    \structure{Computation as proof reduction}
    \begin{itemize}
    \item Functional programming
    \item Types --- preservation and progress properties
    \end{itemize}

  \column{0.55\textwidth}
    \begin{tikzcd}[arrow style=tikz, /tikz/>={To[scale=0.8]}, arrows={very thick, start anchor=real east, end anchor=real west}, /tikz/nodes={inner xsep=0}]
      \begin{tikzpicture}[scale=0.75]
        % \fill [gray] (0:0) -- ++(60:1.45) -- ++(0:-1.45) -- cycle;
        \fill [example text.fg] (-60:-2) -- ++(-60:0.5) -- (60:1.5) -- ++(60:0.5) -- cycle;
        \path (-60:-2.55 -| -60:-2.05) -- ++(2.05,0) -- (60:2.05) -- ++(-2.05, 0) -- cycle;
      \end{tikzpicture}
      \rar[structure] &
      \begin{tikzpicture}[scale=0.75]
        % \fill [gray] (0:0) -- ++(60:1.45) -- ++(0:-1.45) -- cycle;
        \fill [example text.fg] (-60:-2) -- ++(-60:0.5) -- (60:1.5) -- ++(60:0.5) -- cycle;
        \fill [structure] (-60:-2.55) -- ++(-60:0.5) -- (60:2.05) -- ++(60:0.5) -- cycle;
      \end{tikzpicture}
    \end{tikzcd}
    \structure{Computation as proof construction}
    \begin{itemize}
    \item Committed choice bottom-up logic programming
    \item Only minimal types --- can get stuck
    \end{itemize}
  \end{columns}
\end{frame}


\begin{frame}[fragile]
  \frametitle{\emph{Concurrent} proof reduction and proof construction}

  \begin{columns}[b,onlytextwidth]
  \column{0.3\textwidth}
    \begin{center}
      \fcolorbox{structure!30}{structure!10}{$
        \begin{gathered}
          \text{Tasks:} \\
          \begin{tikzcd}[ampersand replacement=\&, row sep=1.5ex, column sep=1.8em, arrow style=tikz, /tikz/>={To[scale=0.8]}, arrows={very thick, start anchor=real east, end anchor=real west}, /tikz/nodes={inner ysep=0}]
            \circtask[structure]{1.7ex}
              \rar[structure] \&
            \sqtask[gray]{1.7ex}
            \\
            \circtask[example text.fg]{1.7ex}
              \rar[example text.fg] \&
            \sqtask[gray]{1.7ex}
          \end{tikzcd}
        \end{gathered}
      $}
    \end{center}

  \column{0.6\textwidth}
    \begin{tikzcd}[row sep=2ex, arrow style=tikz, /tikz/>={To[scale=0.8]}, arrows={very thick}, /tikz/nodes={inner ysep=0}]
      &
      \sqtask[gray]{1.7ex} \, \circtask[example text.fg]{1.7ex}
        \drar[example text.fg, start anchor=east, end anchor=north west] &
      \\
      \circtask[structure]{1.7ex} \, \circtask[example text.fg]{1.7ex}
        \urar[structure, start anchor=north east, end anchor=west]
        \drar[example text.fg, start anchor=south east, end anchor={[yshift=-axis_height]north west}]
        \ar[rr, alert, dashed, dash phase=0.5em, start anchor=real east, end anchor=real west] &
      &
      \sqtask[gray]{1.7ex} \, \sqtask[gray]{1.7ex}
      \\
      &
      \circtask[structure]{1.7ex} \, \sqtask[gray]{1.7ex}
        \urar[structure, start anchor={[yshift=-axis_height]north east}, end anchor=south west] &
    \end{tikzcd}
  \end{columns}

  \vspace*{3ex}

  \setbeamerfont*{itemize/enumerate body}{size=\small}
  \structure{Concurrent equality:} Interleavings of independent tasks are indistinguishable,
  so that they appear to be truly concurrent.
  \begin{itemize}
  \item This diagram commutes.
  \end{itemize}

  \vspace*{1ex}

  \begin{columns}[t,onlytextwidth]
  \column{0.425\textwidth}
    \structure{Concurrent proof reduction}
    \begin{tikzcd}[row sep=0, every cell/.append style={rectangle, inner sep=0}, arrow style=tikz, /tikz/>={To[scale=0.8]}, arrows={very thick, start anchor=east, end anchor=west}]
      & \drar[example text.fg, end anchor={[yshift=0.75ex]west}] & \\
      \begin{tikzpicture}[scale=0.5]
        \fill [gray] (0:0) -- ++(60:2) -- ++(0:-2) -- cycle;
        \clip (-60:-2.55) -- ++(-60:0.5) -- (60:2.05) -- ++(60:0.5) -- cycle;
        \fill [structure] (-60:-1.35) -- ++(60:1.2) -- ++(0:-1.2) -- cycle;
        \fill [example text.fg] (60:1.35) -- ++(60:1.2) -- ++(0:-1.2) -- cycle;
      \end{tikzpicture}
      \urar[structure, start anchor={[yshift=0.75ex]east}] \drar[example text.fg, start anchor={[yshift=-0.75ex]east}] \ar[rr, alert, dashed] & &
      \begin{tikzpicture}[scale=0.5]
        \fill [gray] (0:0) -- ++(60:2) -- ++(0:-2) -- cycle;
        \clip (-60:-2.55) -- ++(-60:0.5) -- (60:2.05) -- ++(60:0.5) -- cycle;
        % \fill [example text.fg] (60:1.35) -- ++(60:1.2) -- ++(0:-1.2) -- cycle;
      \end{tikzpicture}
      \\
      & \urar[structure, end anchor={[yshift=-0.75ex]west}] &
    \end{tikzcd}
    % \begin{tikzcd}[every cell/.append style={rectangle, inner xsep=0.15em, inner ysep=0.5em, draw}]
    %   & \drar[example text.fg] & \\
    %   {}
    %     \urar[structure]
    %     \drar[example text.fg]
    %     \ar[rr, alert, dashed] & &
    %   {}
    %   \\
    %   & \urar[structure] &
    % \end{tikzcd}

    \begin{itemize}
    \item SILL: Proofs as processes
    \end{itemize}

    \showcolumnwidth

  \column{0.5\textwidth}
    \structure{Concurrent proof construction}
    % \begin{tikzcd}[/tikz/nodes={draw}]
    %   & \drar[example text.fg] & \\
    %   {}
    %     \urar[structure]
    %     \drar[example text.fg]
    %     \ar[rr, alert, dashed] & &
    %   {}
    %   \\
    %   & \urar[structure] &
    % \end{tikzcd}
    %
    \begin{tikzcd}[row sep=0, every cell/.append style={rectangle, inner sep=0}, arrow style=tikz, /tikz/>={To[scale=0.8]}, arrows={very thick, start anchor=east, end anchor=west}]
      & \drar[example text.fg, end anchor={[yshift=0.75ex]west}] & \\
      %
      \begin{tikzpicture}[scale=0.5, baseline={(0:0 |- 60:1)}]
        \fill [gray] (0:0) -- ++(60:1.45) -- ++(0:-1.45) -- cycle;
        \path (0:0) -- ++(60:2) -- (60:2 |- 60:2.55) -- ++(0:-2) -- (-60:-2) -- cycle;
      \end{tikzpicture}
      \urar[structure, start anchor={[yshift=0.75ex]east}] \drar[example text.fg, start anchor={[yshift=-0.75ex]east}] \ar[rr, alert, dashed]
      &&
      \begin{tikzpicture}[scale=0.5, baseline={(0:0 |- 60:1)}]
        \fill [gray] (0:0) -- ++(60:1.45) -- ++(0:-1.45) -- cycle;
        \begin{scope}
          \path [preaction={fill=structure}, clip] (-60:-2) -- ++(-60:0.5) -- (60:1.5) -- ++(60:0.5) -- cycle;
          \fill [example text.fg] foreach \x in {0,0.4,0.8,1.2,1.6} { ($(-60:-2)+(\x,0)$) -- ++(0.2,0) -- ++(-60:0.5) -- ++(-0.2,0) -- cycle };
        \end{scope}
        \begin{scope}
          \path [preaction={fill=example text.fg}, clip] (-60:-2.55) -- ++(-60:0.5) -- (60:2.05) -- ++(60:0.5) -- cycle;
          \fill [structure] foreach \x in {0,0.4,0.8,1.2,1.6,2.0,2.4} { ($(-60:-2.55)+(\x,0)$) -- ++(0.2,0) -- ++(-60:0.5) -- ++(-0.2,0) -- cycle };
        \end{scope}
      \end{tikzpicture}
      % \equiv
      % \begin{tikzpicture}[scale=0.5, baseline={(0:0 |- 60:1)}]
      %   \fill [gray] (0:0) -- ++(60:1.45) -- ++(0:-1.45) -- cycle;
      %   \fill [example text.fg] (-60:-2) -- ++(-60:0.5) -- (60:1.5) -- ++(60:0.5) -- cycle;
      %   \fill [structure] (-60:-2.55) -- ++(-60:0.5) -- (60:2.05) -- ++(60:0.5) -- cycle;
      % \end{tikzpicture}
      \\
      & \urar[structure, end anchor={[yshift=-0.75ex]west}] &
    \end{tikzcd}
    
    \begin{itemize}
    \item CLF: Proofs as multiset rewriting traces
    \end{itemize}

    \showcolumnwidth
  \end{columns}
\end{frame}


\begin{frame}[plain,fragile]
  \onslide<+->

  \begin{center}
    \begin{tikzcd}[column sep=6em, arrow style=tikz, /tikz/>={To[scale=0.8]}, arrows={very thick, start anchor=real east, end anchor=real west}]
      \begin{tikzpicture}[scale=0.75]
        \fill [gray] (0:0) -- ++(60:2) -- ++(0:-2) -- cycle;
        \clip (-60:-2.55) -- ++(-60:0.5) -- (60:2.05) -- ++(60:0.5) -- cycle;
        \fill [structure] (-60:-1.35) -- ++(60:1.2) -- ++(0:-1.2) -- cycle;
        \fill [example text.fg] (60:1.35) -- ++(60:1.2) -- ++(0:-1.2) -- cycle;
      \end{tikzpicture}
      \rar[<->, "\text{\normalsize\alt<-.>{?}{session types}}"] &
      \begin{tikzpicture}[scale=0.75]
        \fill [gray] (0:0) -- ++(60:1.45) -- ++(0:-1.45) -- cycle;
        \fill [example text.fg] (-60:-2) -- ++(-60:0.5) -- (60:1.5) -- ++(60:0.5) -- cycle;
        \fill [structure] (-60:-2.55) -- ++(-60:0.5) -- (60:2.05) -- ++(60:0.5) -- cycle;
      \end{tikzpicture}
    \end{tikzcd}
  \end{center}

  \begin{center}
    \begin{tikzpicture}
      \graph [grow right sep] {
        sr / "Multiset\\rewriting" [align=flush center]
         -> [bend left]
        c / "Choreography"
         -> [bend left]
        p / "Message-passing\\process" [align=flush center];

        p -> [bend left] c -> [bend left] sr;
      };

      \coordinate (sr-c) at ($(sr.east)!0.5!(c.west)$);
      \draw [dashed] ($(sr-c)!3em!90:(c)$) -- ($(sr-c)!3em!-90:(c)$);
    \end{tikzpicture}
  \end{center}

  \begin{block}<+->{Thesis statement}
    Session types bridge the different notions of concurrency that arise in proof-reduction and proof-construction interpretations of linear logic.
    \begin{itemize}
    \item Guarantees of session types thus apply to (a fragment of) logic programming!
    \end{itemize}
  \end{block}
\end{frame}


\begin{frame}
  \frametitle{Representing natural numbers as binary strings}

  \begin{itemize}
  \item Strings of $\bit{0}$s and $\bit{1}$s beginning with an $\eps$.
    \begin{equation*}
      \!\begin{aligned}
        \val{\eps \strcons \bit{1}} &= 1 \\
        \val{\eps \strcons \bit{0} \strcons \bit{1}} &= 1 \\
        % \val{\eps \strcons \bit{1} \strcons \bit{0}} &= 2 \\
        \val{\eps \strcons \bit{1} \strcons \bit{0} \strcons \bit{1} \strcons \bit{0} \strcons \bit{1} \strcons \bit{0}} &= 42
      \end{aligned}
    \end{equation*}

  \item $\cntr{n}$ is the canonical representation of $n$: no leading $\bit{0}$s.
    % Canonical representative of $n$ is string $\cntr{n}$ with no leading $\bit{0}$s:
    \begin{equation*}
      \!\begin{aligned}
        \cntr{0} &= \eps \\
        \cntr{2n} &= \cntr{n} \strcons \bit{0} \text{, if $2n > 0$} \\
        \mathllap{\cntr{2n+1}} &= \cntr{n} \strcons \bit{1}
      \end{aligned}
    \end{equation*}
  \end{itemize}
\end{frame}


\begin{frame}[fragile]
  \frametitle{Computing successors via string rewriting}

  \begin{itemize}
  \item $\inc$ means successor.   Grade school arithmetic as string rewriting:
    % \begin{block}{Adequacy}
      \begin{equation*}
        \cntr{n} \strcons \inc \trans* \cntr{n+1}
      \end{equation*}  
    % \end{block}

  % \item Grade school arithmetic as string rewriting:
    \begin{minipage}{\linewidth}
    \begin{columns}[T,totalwidth=\linewidth]
    \column{0.5\linewidth}
      \fcolorbox{structure!30}{structure!10}{%
        \begin{minipage}{0.795\linewidth}
          \begin{description}
          \item[Ends in $\bit{1}$?] $\bit{1} \strcons \inc \trans \inc \strcons \bit{0}$
          \item[Is $\eps$?] $\eps \strcons \inc \trans \eps \strcons \bit{1}$
          \item[Ends in $\bit{0}$?] $\bit{0} \strcons \inc \trans \bit{1}$
          \end{description}
        \end{minipage}
      }

      \showcolumnwidth

    \column{0.5\linewidth}
      $\cntr{1} \strcons \inc \trans* \cntr{2}$\\
      % \begin{equation*}
        $\eps \strcons \bit{1} \strcons \inc
          \trans
        \eps \strcons \inc \strcons \bit{0}
          \trans
        \eps \strcons \bit{1} \strcons \bit{0}$
      % \end{equation*}  

      \showcolumnwidth

    \end{columns}
  \end{minipage}
  % \fcolorbox{structure!30}{structure!10}{$
  %   \!\begin{gathered}
  %     \text{Grade school arithmetic:}\\
  %     \!\begin{aligned}
  %       \bit{1} \strcons \inc
  %         &\trans \inc \strcons \bit{0} \\
  %       % 
  %       \eps \strcons \inc
  %         &\trans \eps \strcons \bit{1} \\
  %       % 
  %       \bit{0} \strcons \inc
  %         &\trans \bit{1}
  %     \end{aligned}
  %   \end{gathered}
  % $}

  \item Several \emph{concurrent} successor instructions at once:
    \begin{equation*}
      \begin{tikzcd}[every cell/.append style={shape=rectangle, inner sep=0.2em}]
        & & \drar & \\
        |[asymmetrical rectangle]|
        \eps \strcons \bit{1} \strcons \inc \strcons \inc
          \rar & |[asymmetrical rectangle]| \urar \drar \ar[rr,dashed] & & {} \\
        & & \urar &
      \end{tikzcd}
    \end{equation*}
  \end{itemize}
\end{frame}


\begin{frame}
  Repeat circle diagram

  Where are the processes hiding?
\end{frame}


\begin{frame}
  \frametitle{Letters as processes}

  \begin{equation*}
    \eps \strcons \bit{1} \strcons \inc \strcons \inc \text{\quad picture}
  \end{equation*}
  \begin{itemize}
  \item Letters as processes (cf.\ [Miller '92])
  \item Strings as process chains
  \item Rewriting rules as process transformations/specifications
  \end{itemize}

  This is not a local implementation; it assumes an omniscient coordinator.
  \begin{itemize}
  \item Iris picture here.
  \item Above is closer to a one-dimensional cellular automaton.
  \end{itemize}
\end{frame}


\begin{frame}
  \frametitle{Letters as messages \emph{and} processes}

  \begin{equation*}
    \eps \strcons \bit{1} \strcons \inc \strcons \inc \text{\quad picture}
  \end{equation*}
  \begin{itemize}
  \item Some letters as processes and others as messages
  \item Strings as chains of processes and messages
  \item Rewriting rules as \emph{locally}-executing process implementations
  \end{itemize}

  Assignment of letters to process/message roles is a \emph{\textcolor{structure}{choreography}}
\end{frame}


\begin{frame}[plain]
  Have we done anything but draw pictures?

  Examples of ill-formed choreographies
  \begin{itemize}
  \item Two processes
  \item Two messages
  \item Two messages received -- doesn't happen atomically in standard process calculi
  \item Message sent in premise; message received in conclusion
  \end{itemize}

  Locality: premise has one process and at most one message received, and conclusion has no messages received.

  Bisimulation condition?
\end{frame}


\begin{frame}
  repeat rewriting--choreography--process picture
\end{frame}


\begin{frame}[fragile]
  \frametitle{A picture is worth a few lines of code}

  \begin{columns}
  \column{0.5\textwidth}
  Picture of process reduction

  \column{0.5\textwidth}
    \begin{Verbatim}[gobble=4, commandchars=\\\{\}]
    proc \textcolor{structure}{bit1} =
      case (recvR ()) of
        \textcolor{exampletext}{s} => sendL \textcolor{exampletext}{s}; \textcolor{structure}{bit0}
    \end{Verbatim}

    \begin{align*}
      % \MoveEqLeft[0.5]
      &\exec{\text{\Verb[commandchars=\\\{\}]+\textcolor{structure}{bit1}+}} \mskip\thickmuskip \msg{\text{\Verb[commandchars=\\\{\}]+\textcolor{exampletext}{s}+}} %\\[-0.5ex]
        \trans {\exec{(\text{\Verb[commandchars=\\\{\}]+sendL \textcolor{exampletext}{s}; \textcolor{structure}{bit0}+})}} \\%[1ex]
      %
      % \MoveEqLeft[0.5]
      &\exec{(\text{\Verb[commandchars=\\\{\}]+sendL \textcolor{exampletext}{s}; \textcolor{structure}{bit0}+})} %\\[-0.5ex]
        \trans {\msg{\text{\Verb[commandchars=\\\{\}]+\textcolor{exampletext}{s}+}} \mskip\thickmuskip \exec{\text{\Verb[commandchars=\\\{\}]+\textcolor{structure}{bit0}+}}}
    \end{align*}
  \end{columns}

  Process simulates choreography
\end{frame}


\begin{frame}
  
\end{frame}


\begin{frame}[fragile]
  \frametitle{Sometimes several choreographies are possible}

  Picture of $\eps \strcons \bit{1} \strcons \inc \strcons \inc$ under bits-as-messages choreography

  \begin{verbatim}
proc succ =
{ caseL of
    e => selectR 1; selectL e; closeR
  | 0 => selectR 1; <->
  | 1 => selectR 0; succ }
  \end{verbatim}
\end{frame}

\begin{frame}[fragile]
  \tiny
  \begin{verbatim}
stype Bit = &{ s: Bit , =: Bit , <: bool => Bit' , +: bool => Bit' }
  and Bit' = +{ >: bool /\ Bit }

proc eps :: |- Bit =
{ caseR of
    s => ...
  | + => x <- recvR; selectR >; sendR x; eps -- + }

proc bit b :: Bit |- Bit =
{ caseR of
    s => ...
  | = => selectL <; sendL b; ='
  | < => x <- recvR; selectL <; sendL x; bit' b
  | + => x <- recvR;
         (case b of
            0 => selectR >; sendR x; +
          | 1 => (case x of
                    0 => selectR >; sendR 1; +
                  | 1 => selectR >; sendR 0; selectL s; +)) }

proc bit' b :: Bit' |- Bit' =
{ caseL of
    > => x <- recvL; selectR >; sendR x; bit b }

proc =' :: Bit' |- Bit =
{ caseL of
    > => x <- recvL; selectL =; bit x }

proc + :: Bit |- Bit =
{ caseR of
    < => x <- recvR; selectL +; sendL x; <->
  | eq => <-> }
  \end{verbatim}
\end{frame}

\end{document}

% e 1 + 1 1 =
% e 1 + 1 <1 ='
% e 1 + <1 1' ='
% e 1 +1 1' ='
% e s + 0> 1' ='
% e 1 + 1 0> ='
% e 1 + 1 0> ='
% e 1 + 1 = 0
% e 1 + <1 =' 0
% e 1 +1 =' 0
% e s + 0> =' 0
% e 1 + = 0 0
% e 1 0 0

% e + 1 =
% e +1 ='
% e + 1> =' 

% 0 = -> <0 ='
% b <0 -> <0 b'
% + <0 -> +0
% + <1 -> +1
% b +0 -> b>
% 0 +1 -> 1>
% 1 +1 -> s 0>
% b> 0' -> 0 b>
% b> 1' -> 1 b>
% b> =' -> = b
% + = -> 
% e +0 -> e 0>
% e +1 -> e 1>




% e B 1 + B a1 =  ->*  e B 1 + b1 B =  ->*  e B + b0 B =  ->* 

% 0 =  ->  <0
% 1 =  ->  <1
% b <0  ->  <0 <b
% b <1  ->  <1 <b
% + <0  ->  +0
% 

% B + =
% B 1 +0 B'  ->  B + B' = 1







\begin{frame}
  \frametitle{Title}

  \begin{columns}[t,onlytextwidth]
  \column{0.5\textwidth}
    Concurrency as proof search
    \begin{itemize}
    \item CLF 
    \end{itemize}

  \column{0.5\textwidth}
    Concurrency as proof reduction
    \begin{itemize}
    \item Session-typed $\pi$-calculus
    \end{itemize}
  \end{columns}

  Circle diagram
  \begin{center}
    \begin{tikzpicture}
      \graph [grow right sep=8em, nodes={draw,circle}] {
        caps / "Concurrency\\as\\proof search" [align=flush center]
        -> [bend left]
        "Concurrency\\as\\proof reduction" [align=flush center]
        -> [bend left]
        (caps);
      };
    \end{tikzpicture}
  \end{center}
\end{frame}

\begin{frame}[plain]
  \begin{block}{Thesis statement}
    Session types form a bridge between different notions of concurrency that arise in computational interpretations of linear logic: computation-as-proof-search and computation-as-proof-reduction.
  \end{block}
\end{frame}

\end{document}